%! AP Statistics — Unit 1, Lesson 1.5 — Guided Notes
\documentclass[10pt]{article}
\usepackage{apstats-article}
\usepackage{apstats-boxes}

%=============================================================================
\begin{document}

\pageheader{Unit 1, Lesson 1.5}{Guided Notes}
\namedateperiod

\vspace{0.12in}

% ── OBJECTIVE ────────────────────────────────────────────────────────────────
\begin{objectivebox}
\small By the end of this lesson, I will be able to\dots\\[4pt]
\writeline\\[1pt]
\writeline\\[1pt]
\writeline
\end{objectivebox}

\vspace{0.1in}

% ── VOCABULARY ───────────────────────────────────────────────────────────────
\begin{vocabbox}
\begin{multicols}{2}
\termblanklong{Dotplot}
\termblanklong{Stemplot (stem-and-leaf plot)}
\columnbreak
\termblanklong{Histogram}
\termblanklong{Bin (class interval)}
\termblanklong{Gap / Cluster}
\end{multicols}
\end{vocabbox}

\vspace{0.1in}

% ── HOOK ─────────────────────────────────────────────────────────────────────
\begin{hookbox}
\small\textbf{Three displays, one dataset --- which do you prefer?}

\vspace{6pt}
Ages (years) of 15 summer camp participants:
\[10,\;10,\;10,\;10,\;10,\;11,\;11,\;11,\;11,\;12,\;12,\;12,\;13,\;13,\;14\]

\vspace{4pt}
\renewcommand{\arraystretch}{1.3}
\begin{tabularx}{\linewidth}{@{} >{\bfseries}p{0.17\linewidth}
                                   X X @{}}
\rowcolor{sky}
\textbf{Display} & \textbf{What it shows} & \textbf{Example (ages above)} \\
\midrule
Dotplot   & Every individual value on a number line
          & Dots stacked above 10, 11, 12, 13, 14 \\
Stemplot  & Stem = tens digit; leaf = units digit; values sorted
          & Stem 1 $|$ leaves: 0 0 0 0 0 1 1 1 1 2 2 2 3 3 4 \\
Histogram & Counts grouped into bins; bars touch
          & Bins: {[}10,12{)}: 9 \quad {[}12,14{)}: 5 \quad {[}14,16{)}: 1 \\
\end{tabularx}

\vspace{8pt}
What does each display make easy to see? What does each hide? \blank{5cm}
\end{hookbox}

\vspace{0.1in}

% ── SECTION 1: DOTPLOTS ───────────────────────────────────────────────────────
\begin{notesbox}{Display 1 --- Dotplots}
\small
\begin{multicols}{2}

\textbf{\textcolor{navy}{How to read a dotplot:}}
\begin{itemize}[itemsep=3pt]
  \item The horizontal axis shows the \blank{2.5cm} of the variable.
  \item Each dot represents \blank{2cm}.
  \item The height of a stack of dots shows the \blank{2.2cm} at that value.
  \item Look for: center, spread, shape, \blank{2cm}, \blank{2cm}.
\end{itemize}

\columnbreak

\textbf{\textcolor{navy}{Best for:}}
\begin{itemize}[itemsep=3pt]
  \item \blank{1.5cm} datasets (roughly $n \le $ \blank{0.9cm})
  \item When you want to see \blank{2.5cm} data values
  \item Identifying gaps, clusters, and outliers quickly
\end{itemize}

\vspace{6pt}
\textbf{\textcolor{navy}{Limitation:}}\\
Impractical when $n$ is \blank{2cm} --- dots pile up and overlap.

\end{multicols}
\end{notesbox}

\vspace{0.1in}

% ── SECTION 2: STEMPLOTS ─────────────────────────────────────────────────────
\begin{notesbox}{Display 2 --- Stemplots (Stem-and-Leaf Plots)}
\small
\begin{multicols}{2}

\textbf{\textcolor{navy}{Structure:}}
\begin{itemize}[itemsep=3pt]
  \item \textbf{Stem} = \blank{3.5cm} digits
  \item \textbf{Leaf} = \blank{3.5cm} digit
  \item Leaves are written in \blank{2cm} order.
  \item Always include a \blank{2cm}: e.g., \textbf{7 $|$ 4} = 74.
\end{itemize}

\vspace{6pt}
\textbf{\textcolor{navy}{Complete this stemplot}} for the quiz scores:\\
54, 58, 61, 65, 67, 72, 78, 79, 85

\vspace{4pt}
\begin{tabular}{r|l}
\textbf{Stem} & \textbf{Leaf} \\
\hline
5 & \blank{3.5cm} \\
6 & \blank{3.5cm} \\
7 & \blank{3.5cm} \\
8 & \blank{3.5cm} \\
\end{tabular}

\columnbreak

\textbf{\textcolor{navy}{Best for:}}
\begin{itemize}[itemsep=3pt]
  \item Moderately sized datasets ($n \approx$ \blank{1.8cm})
  \item When you want to \blank{3cm} the actual data values
  \item Comparing two groups with a \blank{2.5cm} stemplot
\end{itemize}

\vspace{6pt}
\textbf{\textcolor{navy}{Advantage over dotplot:}}\\
Data is automatically \blank{2.5cm}, making it easy to find the median.

\vspace{6pt}
\textbf{\textcolor{navy}{Limitation:}}\\
Difficult with very large data values or non-integer data.

\end{multicols}
\end{notesbox}

\newpage

% ── SECTION 3: HISTOGRAMS ─────────────────────────────────────────────────────
\begin{notesbox}{Display 3 --- Histograms}
\small
\begin{multicols}{2}

\textbf{\textcolor{navy}{Key features:}}
\begin{itemize}[itemsep=3pt]
  \item Data grouped into equal-width \blank{2cm} (bins).
  \item Bar height $=$ \blank{2.5cm} (or relative frequency).
  \item Bars \blank{1.5cm} --- no gap between adjacent bins.
  \item Horizontal axis: the \blank{2cm} being measured.
  \item Vertical axis must start at \blank{0.8cm}.
\end{itemize}

\vspace{6pt}
\textbf{\textcolor{navy}{Reading a histogram:}}\\
A histogram shows the \blank{3cm} of the distribution --- but NOT individual values.

\columnbreak

\textbf{\textcolor{navy}{Histogram vs.\ bar chart:}}

\vspace{4pt}
\renewcommand{\arraystretch}{1.5}
\begin{tabularx}{0.46\textwidth}{@{} >{\bfseries}p{0.18\textwidth} X @{}}
\rowcolor{sky}
 & \textbf{Key difference} \\
\hline
Bars touch & \blank{2.8cm} data \\
Bars separate & \blank{2.8cm} data \\
Order matters & \blank{2.8cm} \\
Order arbitrary & \blank{2.8cm} \\
\end{tabularx}

\vspace{6pt}
\textbf{\textcolor{navy}{Best for:}}\\
Large datasets ($n > $ \blank{0.8cm}) where individual values matter less than overall \blank{2cm}.

\end{multicols}
\end{notesbox}

\vspace{0.1in}

% ── WORKED EXAMPLE ───────────────────────────────────────────────────────────
\begin{notesbox}{Worked Example --- Interpreting a Histogram}
\small

\noindent The frequency table below describes the time (minutes) 40 students spent
reading on a Tuesday. Use it to answer the questions.

\vspace{8pt}
\renewcommand{\arraystretch}{1.6}
\begin{tabularx}{\linewidth}{@{} C{0.22\linewidth} C{0.18\linewidth} C{0.22\linewidth} @{}}
\rowcolor{sky}
\textbf{Time (min)} & \textbf{Frequency} & \textbf{Rel.\ Frequency} \\
\midrule
$[0, 20)$   & 4  & \blank{2.5cm} \\
$[20, 40)$  & 12 & \blank{2.5cm} \\
$[40, 60)$  & 15 & \blank{2.5cm} \\
$[60, 80)$  & 6  & \blank{2.5cm} \\
$[80, 100)$ & 3  & \blank{2.5cm} \\
\midrule
\rowcolor{sky}\textbf{Total} & \blank{2cm} & \blank{2.5cm} \\
\end{tabularx}

\vspace{10pt}
\begin{enumerate}[leftmargin=1.5em, itemsep=5pt]
  \item Which bin has the greatest frequency? \blank{4cm}
  \item What percentage of students read for at least 40 minutes?
        Show your work.\quad\blank{5cm}
  \item Describe the shape of this distribution in one word. \blank{3cm}
  \item Write one AP-style sentence interpreting the histogram in context.\\[5pt]
        \writeline\\[1pt]\writeline
\end{enumerate}

\end{notesbox}

\vspace{0.1in}

% ── GUIDED PRACTICE ──────────────────────────────────────────────────────────
\begin{practicebox}
\small
\textbf{Choosing the Right Display.}\\[4pt]
For each situation, identify the \textit{best} display (dotplot, stemplot, or
histogram) and explain your reasoning in one sentence.

\vspace{8pt}
\begin{enumerate}[leftmargin=1.5em, itemsep=8pt]
  \item Heights (cm) of 500 adults in a medical study.\\
        Display: \blank{2.5cm} \quad Reason: \blank{7cm}

  \item Finishing times (seconds) for 8 runners in a 100 m race.\\
        Display: \blank{2.5cm} \quad Reason: \blank{7cm}

  \item Test scores (out of 100) for two AP Statistics classes of 25 students each.\\
        Display: \blank{2.5cm} \quad Reason: \blank{7cm}
\end{enumerate}
\end{practicebox}

\vspace{0.1in}

% ── SPIRAL ───────────────────────────────────────────────────────────────────
\begin{spiralbox}
\small
\begin{multicols}{2}

\textbf{This lesson connects forward to\dots}
\begin{itemize}[itemsep=2pt]
  \item \textbf{Lesson 1.6:} Use today's graphs to describe shape, center, spread,
        and outliers with the SOCS framework.
  \item \textbf{Lesson 1.8:} Boxplots are another quantitative display --- compare
        them to histograms.
  \item \textbf{Unit 4:} Histogram shape connects to probability distributions.
\end{itemize}

\columnbreak

\textbf{Big Idea --- Write in your own words:}\\
\textit{Why does a histogram lose information that a dotplot or stemplot preserves?
When is that trade-off worth making?}\\[2pt]
\writeline\\[1pt]\writeline\\[1pt]\writeline

\end{multicols}
\end{spiralbox}

\vspace{0.1in}

\begin{notesbox}{Additional Notes}
\vspace{0pt}
\writeline\\\writeline\\\writeline\\\writeline\\\writeline\\\writeline
\end{notesbox}

\end{document}
